\section{Exterior packet costs: exact cancellation and retained local jets}
\label{sec:alternating-packet-cost}

The purpose of this calculation is to determine exactly what the alternating
part of the arithmetic tensor construction retains.  A determinant remembers
the sum of the spectral points, whereas proper exterior powers remember all
subset sums.  We compute the difference, including the original arithmetic
metric and every local nilpotent, rather than identify these two observations.
The result is a companion to the symmetric calculation in the
\href{https://doi.org/10.5281/zenodo.22728704}{216-page Split-Zero continuation},
sections ``Coherent packets, joint tensor layers, and their exact spectral
cost'' (CT), ``Density of the original joint boundaries and determinant
dissipation'' (JD), and ``The signed symmetric summand and its complete
spectral cost'' (SS).  Those sections supply the actual source construction;
the exterior comparisons below are proved here.

The standard tensor signs and exterior quotient are those of the
\href{https://stacks.math.columbia.edu/tag/0GWN}{Stacks Project, Tag 0GWN}
and \href{https://stacks.math.columbia.edu/tag/01CF}{Tag 01CF}.
The ordinary permutation convention also agrees with B\"urgisser--Ikenmeyer,
\emph{Explicit Lower Bounds via Geometric Complexity Theory},
\href{https://arxiv.org/abs/1210.8368}{arXiv:1210.8368}, subsection
``A Consequence of Schur-Weyl Duality.''  Its use here is the tensor-action
dictionary, not a complexity lower bound or an RH conclusion.

\subsection{The original source and the unsigned chain average}

Let \(h(t)=\prod_\rho(t-\rho)^{m_\rho}\) be a finite complete packet
of zeros of \(g=2\xi\), with \(d=\deg h\ge1\), and retain
\(E=\C[t]/h\), \(A=M_t\), and all the indicated multiplicities.
For \(1\le k\le d\) and \(M\ge k(d-1)\), the joint source of CT is
\[
 \begin{gathered}
 D_j=-x_j\partial_{x_j},\quad s_j=\tfrac12+it_j,\quad
 v_h=g/h,\quad d\nu_h(t)=|v_h(\tfrac12+it)|^2\frac{dt}{2\pi},\\
 \mathcal T_h^{(k)}P=P(D_1,\ldots,D_k)F_h^{\otimes k},\quad
 \mathcal MF_h=g/h,\quad h(D)F_h=f_0,\\
 \mathcal I_M=\C[s_1,\ldots,s_k]_{\le M}
                  \cap(h(s_1),\ldots,h(s_k)),\quad
 \mathcal B_M=\mathcal T_h^{(k)}\mathcal I_M
                  \subset\mathcal H_k=L^2(d^kx).
 \end{gathered}\tag{EA.1}
\]
The full jet \(J^{(k)}\) is remainder modulo each \(h(s_j)\),
followed by multiplication by \(\upsilon_h^{\otimes k}\), where
\(\upsilon_h=[g/h]_h\) is the original invertible arithmetic unit.
Its inverse remains in the section
\[
 s_hu=\mathcal T_h\operatorname{rem}_h(\upsilon_h^{-1}u).
\]
With all inner products inherited from \(\mathcal H_k\), define
\[
 \begin{gathered}
 A^{(k)}=\sum_{j=1}^k A_j,\quad
 R_M=(1-P_{\mathcal B_M})s_h^{\otimes k},\quad
 G_M=R_M^*R_M,\quad J^{(k)}R_M=1,\\
 \mathcal E_M=\mathcal B_{M+1}\cap\mathcal B_M^\perp,\quad
 B_M=D^{(k)}R_M-R_MA^{(k)},\quad
 Y_M=P_{\mathcal E_M}R_M,\quad C_M=(1-P_{\mathcal B_M})B_M.
 \end{gathered}\tag{EA.2}
\]
Both last maps take values in \(\mathcal E_M\).  These are the
constructed minimum representatives and source maps of CT.12--16, with
\[
 \begin{split}
 G_M&>0,\qquad G_M-G_{M+1}=Y_M^*Y_M,\\
 W_M&=(A^{(k)})^*G_M+G_MA^{(k)}-kG_M
                         =-(Y_M^*C_M+C_M^*Y_M).
 \end{split}\tag{EA.3}
\]
In particular the coefficient \(k\) comes from
\((D^{(k)})^*+D^{(k)}=k\) in the original measure, not from a new scale.

For clarity about the cohomological target, retain the original complex
\(\mathcal C=[V\xrightarrow{\Theta}\mathscr B]\) in degrees \(0,1\),
its quotient \(Q\), and the continuous retraction of SS.1--6:
\[
 \Lambda\Theta=1,\quad s[F]=(1-\Theta\Lambda)F,\quad qs=1,
 \qquad P^0=0,\quad P^1=sq,\quad H^1=\Lambda,
 \quad dH+Hd=1-P.
 \tag{EA.4}
\]
Here \(\Lambda=P_V(1-A_\alpha B_\beta)^{-1}Y_{\alpha,\beta}\)
retains the source's cutoffs and moment projection.  All completions below
are completed projective tensors of those strong-Schwartz locally convex
spaces.  For homogeneous inputs the chain permutation is
\[
 T_\pi(x_1\otimes\cdots\otimes x_k)=
 (-1)^{\sum_{i<j,\,\pi(i)>\pi(j)}|x_i||x_j|}
 x_{\pi^{-1}(1)}\otimes\cdots\otimes x_{\pi^{-1}(k)}.
 \tag{EA.5}
\]
Adjacent interchanges square to one; both braid words have sign
\((-1)^{ab+ac+bc}\).  The differential
\(d(x\otimes y)=dx\otimes y+(-1)^{|x|}x\otimes dy\)
commutes with each such interchange by direct substitution.  Hence
\(T\) is a continuous cochain representation.  Its unsigned average is
\[
 \mathsf A_k=\frac1{k!}\sum_{\pi\in S_k}T_\pi,
 \qquad \mathsf A_k^2=\mathsf A_k,
 \qquad \left.\mathsf A_k\right|_{\text{degree }k}=P_-:=\frac1{k!}
                         \sum_\pi\operatorname{sgn}(\pi)P_\pi.
 \tag{EA.6}
\]
There are \(k!\) ordered factorizations of each permutation, proving
idempotence.  The last equation follows because all top-degree factors
have degree one: \(T_\pi=\operatorname{sgn}(\pi)P_\pi\).

The completed cohomology comparison has explicit inverse maps.  Put
\[
 H_k=\sum_{j=1}^k P^{\widehat\otimes(j-1)}\widehat\otimes H
                    \widehat\otimes1^{\widehat\otimes(k-j)},
 \qquad \overline H_k=\frac1{k!}\sum_\pi T_\pi H_kT_\pi^{-1}.
 \tag{EA.7}
\]
The degree \(-1\) map in slot \(j\) has sign
\((-1)^{\sum_{i<j}|x_i|}\).  In \(dH_k+H_kd\), different-slot
terms cancel, while the same-slot terms telescope to \(1-P^{\otimes k}\)
using (EA.4).  Continuity extends the equality from algebraic tensors;
averaging makes the homotopy commute with \(\mathsf A_k\).  Restricting
\(q^{\widehat\otimes k}\) and \(s^{\widehat\otimes k}\) therefore gives
\[
 H^r(\operatorname{im}\mathsf A_k)=
 \begin{cases}\operatorname{im}(P_-:Q^{\widehat\otimes k}
                                      \to Q^{\widehat\otimes k}),&r=k,\\
 0,&r\ne k.
 \end{cases}\tag{EA.8}
\]
The complementary complex has the complementary cohomology; the direct
sum maps are \(x\mapsto(\mathsf A_kx,(1-\mathsf A_k)x)\) and
\((u,v)\mapsto u+v\).  This proves the comparison on the given
completion by its contraction, without requiring exactness of arbitrary
completed tensor products.

For finite traces use \(\Pi_h^0=0\),
\(\Pi_h^1=\iota_hJ_h\), \(\iota_h=s\sigma_h\),
\(J_h\iota_h=1\), \(J_h\Theta=0\) from SS.8.  Thus
\(\Pi_h^{\otimes k}\mathsf A_k\) is the finite cochain projector
of rank \(\binom dk\), with cohomology inclusion and retraction
\(\sigma_h^{\otimes k}\) and \(j_h^{\otimes k}\) restricted to
the alternating images.  The infinite projector itself is not assigned
this rank.  For \(U(a)=e^{(\log a)A}\), \(a>0\), expansion in a
tensor basis contracts one index cycle at a time, giving
\[
 \operatorname{Tr}_{\wedge^k E}(\wedge^k U(a))
 =\frac1{k!}\sum_\pi\operatorname{sgn}(\pi)
               \prod_{c\in\operatorname{Cycles}(\pi)}
                      \operatorname{Tr}(U(a)^{|c|}).
 \tag{EA.9}
\]
The cochain supertrace is \((-1)^k\) times this trace.  Its validity
on the infinite ambient complex uses finite factorization:
\(\operatorname{Tr}_X(LR)=\operatorname{Tr}_F(RL)\) for finite
\(F\), proved by expanding \(LR\) as a sum of coefficient
functionals times image vectors.  No infinite-dimensional trace is used.

At an invariant support fibre the map is exactly
\((\lambda,x)\mapsto(\lambda,Tx)\).  For a general mask the target is
the original orbit join \(c\lambda=\bigvee_\pi\pi\lambda\), with
the specified transports into it.  The two projected amplitudes sum to
that transport; its kernel is precisely the upstream loss.  A zero
amplitude remains the zero at its target support, not external absence.

\subsection{Unscaled coordinates, the full generator, and all CRT factors}

Let \(e_1,\ldots,e_d\) be the original power basis.  For
\(J=(j_1<\cdots<j_k)\), define the literal alternating orbit vector
and its inclusion and extraction by
\[
 a_J=\sum_{\pi\in S_k}\operatorname{sgn}(\pi)
       e_{j_{\pi(1)}}\otimes\cdots\otimes e_{j_{\pi(k)}},\quad
 I_-e_J=a_J,\quad L_-=\frac1{k!}I_-^*.
 \tag{EA.10}
\]
The coordinate tensor basis gives, exactly,
\[
 I_-^*I_-=k!1,\qquad L_-I_-=1,\qquad I_-L_-=P_-.
 \tag{EA.11}
\]
Different increasing tuples have disjoint tensor supports; each orbit has
\(k!\) entries of square one.  On an arbitrary tensor, \(I_-L_-\)
is its signed orbit average (and kills a repeated index), proving the
last identity.  The exterior quotient \(q_\wedge\) satisfies
\[
 q_\wedge(a_J)=k!\,e_{j_1}\wedge\cdots\wedge e_{j_k},\qquad
 (q_\wedge|_{\operatorname{im}P_-})^{-1}
       (e_{j_1}\wedge\cdots\wedge e_{j_k})=a_J/k!.
 \tag{EA.12}
\]
Every summand acquires a second identical sign in the quotient.  These
are inverse maps with their literal factorials, not unit-length orbit
vectors.

The source measure and constraints are permutation invariant.  Applying
an inverse permutation to a feasible minimum preserves its norm and its
constraint; uniqueness gives equivariance of \(R_M\).  Thus \(G_M\)
and \(G_M^{-1}\) commute with \(P_-\).  Set
\[
 \begin{gathered}
 G_M^-=I_-^*G_MI_-,\quad W_M^-=I_-^*W_MI_-,\quad
 A^-=L_-A^{(k)}I_-,\quad A^{(k)}I_-=I_-A^-,\\
 (G_M^-)^{-1}=L_-G_M^{-1}L_-^*.
 \end{gathered}\tag{EA.13}
\]
The product of the last expression with \(I_-^*G_MI_-\) is
\(L_-G_M^{-1}P_-G_MI_-=L_-I_-=1\).  This proves the inverse
compression with both factors \(1/k!\) and its commuting projection.

Every matrix coefficient of \(A^-\) is specified by
\[
 A^-(e_{j_1}\wedge\cdots\wedge e_{j_k})
 =\sum_{r=1}^k\sum_{i=1}^d A_{i j_r}
       e_{j_1}\wedge\cdots\wedge e_i\wedge\cdots\wedge e_{j_k}.
 \tag{EA.14}
\]
A repeated index is zero; otherwise reorder the indices and multiply
by the sign of that reordering.  The diagonal coefficient is
\(\sum_{j\in J}A_{jj}\).  Tensor-slot expansion proves (EA.14);
the common factor \(k!\) in (EA.12) cancels in the intertwining,
so this is also the matrix in the unscaled coordinates (EA.10).

The complete CRT map is reduction modulo each \((t-\rho)^{m_\rho}\),
with inverse the unique polynomial of degree below \(d\) having all
prescribed jets.  It yields
\[
 \begin{gathered}
 \wedge^k E\simeq
 \bigoplus_{\substack{0\le\ell_\rho\le m_\rho\\\sum\ell_\rho=k}}
             \bigotimes_\rho\wedge^{\ell_\rho}
                           (\C[x_\rho]/x_\rho^{m_\rho}),\\
 w_\ell=\prod_\rho\binom{m_\rho}{\ell_\rho},\qquad
 A^-|_\ell=\lambda_\ell+N_\ell,\qquad
 \lambda_\ell=\sum_\rho\ell_\rho\rho.
 \end{gathered}\tag{EA.15}
\]
Fix an order of the distinct roots.  The quotient map sends a tensor
of local wedges to their concatenated global wedge, and its inverse
groups each global basis wedge into those root blocks.  In unscaled
tensor images the map is the signed sum over the
\((\ell_\rho)\)-shuffles of the tensor of local alternating sums.
Each permutation factors uniquely into a shuffle and within-block
permutations; their signs multiply.  Hence that signed shuffle sum is
exactly the global \(a_J\), with no additional multinomial factor.
The capacity bounds and dimensions in (EA.15) follow by choosing each
local subset.  Distinct occupancies remain separate even when their
\(\lambda_\ell\) coincide.

\begin{theorem}[Exact exterior nilpotent index]
For one local block \(\C[x]/x^m\), multiplication by \(x\) induces
on \(\wedge^r\), \(0\le r\le m\), an additive nilpotent of index
\(H+1\), where \(H=r(m-r)\).  Its last nonzero transition is
\[
 (N^{[r]})^H(1\wedge x\wedge\cdots\wedge x^{r-1})
 =H!\prod_{j=0}^{r-1}\frac{j!}{(m-r+j)!}
            x^{m-r}\wedge\cdots\wedge x^{m-1}.
 \tag{EA.16}
\]
An empty wedge is the scalar one.  On sector \(\ell\) of (EA.15)
the index is exactly \(1+H_\ell\), with
\[
 H_\ell=\sum_\rho\ell_\rho(m_\rho-\ell_\rho).
 \tag{EA.17}
\]
\end{theorem}
\begin{proof}
In the basis of increasing exponents, each nonzero application raises
their sum by one.  Its smallest and largest values differ by \(H\),
so the power \(H+1\) vanishes.  The matrix of \(e^{uN}\) has
entry \(u^{a-b}/(a-b)!\) for \(a\ge b\), zero otherwise.
Its exterior action equals \(e^{uN^{[r]}}\): differentiating the wedge
of the \(r\) factors gives (EA.14), with identity initial value;
coefficient recursion uniquely solves this polynomial differential
equation.  Thus its bottom-to-top coefficient is
\[
 u^H\det\left[\frac1{(m-r+i-j)!}\right]_{0\le i,j<r}
 =u^H\prod_{j=0}^{r-1}\frac{j!}{(m-r+j)!}.
 \tag{EA.18}
\]
Negative factorial entries mean zero.  Factor \(1/(m-r+i)!\)
out of row \(i\); the remaining entry is the falling factorial
of \(m-r+i\) of degree \(j\).  These are monic polynomials of
successive degrees.  Column elimination converts their determinant to
the Vandermonde at \(m-r,m-r+1,\ldots,m-1\), whose product is
\(\prod_{j=0}^{r-1}j!\).  This proves (EA.18), including entries
where the falling factorial is zero.  Extracting the coefficient
\(u^H/H!\) gives (EA.16), which is strictly positive.  For
\(r=0,m\), the space is a line, \(H=0\), and the zero
operator has index one as asserted.

The local summands of \(N_\ell\) commute.  Total exponent again
bounds its index by \(H_\ell+1\).  The transition from the tensor
of local bottom wedges to the tensor of top wedges has coefficient
\[
 H_\ell!\prod_\rho\prod_{j=0}^{\ell_\rho-1}
                   \frac{j!}{(m_\rho-\ell_\rho+j)!}>0.
 \tag{EA.19}
\]
Indeed the multinomial coefficient distributes \(H_\ell\)
applications among the local last powers (EA.16).  This proves the
exact index while (EA.14)--(EA.15) retain the complete operator, not
merely that index or its trace.
\end{proof}

\subsection{Subset moments, determinant cancellation, and departure}

List \(\rho_1,\ldots,\rho_d\) with all algebraic multiplicities,
and put \(\delta_i=\operatorname{Re}\rho_i-1/2\),
\(\sigma_1=\sum_i\delta_i\), \(\sigma_2=\sum_i\delta_i^2\),
\(N_-=\binom dk\), \(\lambda_J=\sum_{j\in J}\rho_j\),
\(\delta_J=\sum_{j\in J}\delta_j\).  Triangularizing \(A\)
and ordering increasing basis wedges by index weight triangularizes
(EA.14).  Its diagonal is the complete list of \(\lambda_J\).
Consequently, without discarding any nilpotent,
\[
 \det(z-A^-)=\prod_{|J|=k}(z-\lambda_J),\qquad
 \sum_{k=0}^d q^k\operatorname{Tr}(\wedge^k U(a))
                  =\det(1+qU(a))=\prod_{i=1}^d(1+qa^{\rho_i}).
 \tag{EA.20}
\]
The \(k=0\) coefficient here is the scalar trace one.
For \(d\ge2\) write
\[
 \alpha=\binom{d-1}{k-1},\qquad
 \beta=\binom{d-2}{k-1},\qquad
 \gamma=\binom{d-2}{k-2},
 \tag{EA.21}
\]
where an out-of-range lower index gives zero.  Each index belongs to
\(\alpha\) subsets and each distinct pair to \(\gamma\) subsets.
Since \(\alpha-\gamma=\beta\), direct expansion gives
\[
 \sum_{|J|=k}\delta_J=\alpha\sigma_1,\qquad
 \sum_{|J|=k}\delta_J^2=\beta\sigma_2+\gamma\sigma_1^2,
 \qquad
 \sum_{|J|=k}(\delta_J-k\sigma_1/d)^2
                 =\beta(\sigma_2-\sigma_1^2/d).
 \tag{EA.22}
\]
For the last equality use \(\alpha=N_-k/d\) and
\(\gamma-N_-k^2/d^2=-\beta/d\), immediate from the
factorials.  When \(d=k=1\), the first and second sums are simply
\(\delta_1\) and \(\delta_1^2\), and the centered sum is zero.

In the original Gram define
\[
 \begin{split}
 \eta_M^-&=\operatorname{Tr}((G_M^-)^{-1}(A^-)^*G_M^-A^-)
                                      -\sum_{|J|=k}|\lambda_J|^2,\\
 S_M^-&=(G_M^-)^{-1/2}W_M^-(G_M^-)^{-1/2},\quad
 \epsilon_M^-=\|S_M^-\|_{\rm op},\\
 \mathcal D_k^-&=\sum_{|J|=k}|\delta_J|\\
 &=\sum_\ell w_\ell\left|\sum_\rho\ell_\rho
                                  (\operatorname{Re}\rho-1/2)\right|.
 \end{split}\tag{EA.23}
\]
The square root is solely the isometry computing the displayed
Hermitian pencil; none of \(I_-\), \(R_M\), \(G_M\), \(Y_M\), \(C_M\), or the original
arithmetic unit is replaced or rescaled.

\begin{theorem}[Exact exterior metric moments]
For \(d\ge2\), the constructed source matrices satisfy
\[
 \eta_M^-\ge0,\qquad
 \operatorname{Tr}S_M^-=2\alpha\sigma_1,\qquad
 \operatorname{Tr}((S_M^-)^2)
        =4(\beta\sigma_2+\gamma\sigma_1^2)+2\eta_M^-.
 \tag{EA.24}
\]
For \(d=k=1\), replace the bracket by \(\delta_1^2\) and
the first moment by \(\delta_1\).
\end{theorem}
\begin{proof}
Set \(T=(G_M^-)^{1/2}A^-(G_M^-)^{-1/2}\).  A unitary
triangularization follows inductively by choosing a unit eigenvector,
completing it to an orthonormal basis, and treating the remaining
diagonal block.  In that basis \(T\) has diagonal \(\lambda_J\).
Its strictly upper triangular entries have squared sum \(\eta_M^-\),
proving nonnegativity.  By (EA.3) and (EA.13),
\(S_M^-=T+T^*-k1\).  Its diagonal is \(2\delta_J\);
each strictly upper entry contributes its squared modulus twice to
the squared trace.  Thus that trace is
\(4\sum_J\delta_J^2+2\eta_M^-\); (EA.22) proves the result.
\end{proof}

In degree \(k=d\), (EA.14) reduces to the scalar
\(A^-=\operatorname{Tr}A\).  Therefore
\[
 W_M^-=2\sigma_1G_M^-,\quad \eta_M^-=0,\quad
 \mathcal D_d^-=|\sigma_1|,\quad \epsilon_M^-=2|\sigma_1|.
 \tag{EA.25}
\]
For a reflection-stable packet \(\sigma_1=0\), so the whole
determinant control vanishes even if individual offsets do not.
For example the algebraic list \(1/2+a,1/2-a\), with \(a\in\R\), has degree-one
cost \(2|a|\), but degree-two cost zero; this illustrates the
exact cancellation and does not assert such zeros of \(\xi\).
Every local determinant nilpotent also vanishes by (EA.17).

For every proper degree \(1\le k<d\), instead \(\beta>0\).
Equation (EA.22) shows \(\mathcal D_k^-=0\) exactly when all
\(\delta_i=0\).  Indeed its nonnegative second moment then
vanishes and forces \(\sigma_2=0\), and the converse is immediate.
Moreover if any \(m_\rho>1\), some sector retains a nonzero
nilpotent in every proper degree: choose
\[
 \max(1,k-d+m_\rho)\le\ell_\rho\le\min(m_\rho-1,k).
 \tag{EA.26}
\]
This integer interval is nonempty for \(1\le k<d\); the remaining
degree is between zero and the remaining capacity \(d-m_\rho\),
so it can be allocated among the other blocks.  Equation (EA.17)
is then positive.  That operator is not diagonalizable, hence cannot
be unitarily diagonal for any positive Gram.  The squared sum in
the proof of (EA.24) consequently gives \(\eta_M^->0\).

\subsection{Exact relation-layer cost and the complementary information}

Restrict the source maps, retaining their Hilbert targets:
\[
 \begin{gathered}
 Y_M^-=Y_MI_-,\quad C_M^-=C_MI_-,\quad
 \Delta_M^-=G_M^--G_{M+1}^-=(Y_M^-)^*Y_M^-,\\
 W_M^-=-((Y_M^-)^*C_M^-+(C_M^-)^*Y_M^-).
 \end{gathered}\tag{EA.27}
\]
Their target is the alternating part of the original layer.  Its
dimension is exactly
\[
 r_{\rm lay}^-=p_{\le k}\bigl(M+1-k(k-1)/2\bigr),
 \tag{EA.28}
\]
where \(p_{\le k}(n)\) counts partitions into at most \(k\)
parts, is zero for \(n<0\), and equals one at zero.
To prove this, an alternating polynomial vanishes on every diagonal
\(s_i=s_j\), so it is divisible by each \(s_j-s_i\).  These
distinct irreducibles are relatively prime in the polynomial ring,
so their product
\(V_k(s)=\prod_{i<j}(s_j-s_i)\) divides it.  The quotient is
symmetric because numerator and product have the same permutation
sign.  Division by this literal \(V_k\) and multiplication by it
are inverse vector-space maps, shifting degree by \(k(k-1)/2\).
Symmetric monomial sums of a fixed degree are indexed by partitions.
The alternating full-jet map is onto at the stated cutoff: average a
complete tensor remainder by \(P_-\), retaining the commuting
arithmetic unit.  The two successive kernel dimensions therefore
differ by the new exact-degree dimension (EA.28).  The injective
source map and orthogonal quotient identify this difference with the
actual layer.  No replacement of the source measure is involved.

Write
\[
 r_M^-=\operatorname{rank}Y_M^-,\quad
 \tau_M^-=\operatorname{Tr}((G_M^-)^{-1}\Delta_M^-),\quad
 \chi_M^-=\operatorname{Tr}((G_M^-)^{-1}(C_M^-)^*C_M^-).
 \tag{EA.29}
\]
Then all costs are in the same original source:
\[
 \begin{gathered}
 n_\pm(W_M^-)\le
    \min\{r_M^-,\operatorname{rank}C_M^-,r_{\rm lay}^-\},\\
 \mathcal D_k^-\le\tfrac12\|S_M^-\|_1
                 \le\sqrt{\tau_M^-\chi_M^-},\qquad
 \mathcal D_k^-\le r_M^-\epsilon_M^-.
 \end{gathered}\tag{EA.30}
\]
For the inertia assertion, the quadratic form (EA.27) vanishes on
\(\ker Y_M^-\).  A strictly positive or negative subspace of
larger dimension than \(\operatorname{rank}Y_M^-\) intersects
that kernel, a contradiction.  Apply the same argument to \(C_M^-\).
In the triangular basis of (EA.24), the absolute diagonal sum is
\(2\mathcal D_k^-\); a spectral expansion bounds that sum by
the trace norm.  Finally
\(\|B^*C\|_1\le\|B\|_{\rm HS}\|C\|_{\rm HS}\): choose
a unitary attaining the trace norm and apply Cauchy--Schwarz to the
two Hilbert--Schmidt coefficient arrays.  Apply this inequality to
\(B=Y_M^-(G_M^-)^{-1/2}\), \(C=C_M^-(G_M^-)^{-1/2}\)
and their adjoints in (EA.27).  Their squared norms are exactly
\(\tau_M^-,\chi_M^-\).  At most \(2r_M^-\) nonzero
eigenvalues remain, proving the last bound.

The original reflection \(u(t)\mapsto\overline{u(1-\overline t)}\),
\(F(x)\mapsto x^{-1}\overline{F(1/x)}\), restricted to a
reflection-stable complete packet, commutes with permutations.  Its
tensor version sends \(W_M\) to its negative and preserves the
original Gram by CT.16a--b.  Thus it restricts through (EA.10) and
pairs all nonzero generalized eigenvalues of \((W_M^-,G_M^-)\).
With \(p_M^-=n_+(W_M^-)=n_-(W_M^-)\), one obtains
\[
 \begin{gathered}
 p_M^-\le\min\{r_M^-,\operatorname{rank}C_M^-,r_{\rm lay}^-,
                                      \lfloor N_-/2\rfloor\},\\
 \mathcal D_k^-\le p_M^-\epsilon_M^-,\qquad
 p_M^-(\epsilon_M^-)^2\ge2\beta\sigma_2+\eta_M^-.
 \end{gathered}\tag{EA.31}
\]
The \(2p_M^-\) nonzero eigenvalues have absolute value at most
\(\epsilon_M^-\).  Bound their absolute and squared sums and
use (EA.24) with \(\sigma_1=0\).  No division by \(p_M^-\)
is made.  For \(d=k=1\), (EA.25) gives the zero reflected form
directly, so the binomial convention in (EA.31) is unnecessary.

For a further exact comparison, every eigenvalue \(\lambda_J\)
of \(A^-\) has a nonzero eigenvector \(v\).  Its original
Rayleigh quotient is
\(v^*W_M^-v/(v^*G_M^-v)=2\operatorname{Re}\lambda_J-k=2\delta_J\).
Consequently
\[
 \epsilon_M^-\ge2\max_{|J|=k}|\delta_J|,
 \qquad
 \mathcal D_k^-\ge
 \frac{\beta\sigma_2+\gamma\sigma_1^2}{k\delta_*}
 \quad(\delta_*:=\max_i|\delta_i|>0,\ d\ge2).
 \tag{EA.32}
\]
The second inequality sums
\(\delta_J^2\le k\delta_*|\delta_J|\).
The first proof deliberately uses eigenvectors of the exterior operator
itself; a wedge of repeated original eigenvectors may be zero for a
defective root and is not substituted for \(v\).

The ordered tensor has occupancy \(n\) with multiplicity
\(m_n=k!/\prod_i n_i!\).  The alternating image retains one
copy exactly when \(n\) is squarefree.  Its complementary
characteristic polynomial, trace, and absolute spectral cost therefore
have multiplicities \(m_n-\mathbf1_{\{n_i\le1\ \forall i\}}\).
In particular
\[
 \mathcal D_k=\mathcal D_k^-+
 \sum_{|n|=k}(m_n-\mathbf1_{\{n_i\le1\ \forall i\}})|n\cdot\delta|,
 \qquad
 \mathcal D_k^s-\mathcal D_k^-=
             \sum_{\substack{|n|=k\\\exists i:n_i\ge2}}|n\cdot\delta|.
 \tag{EA.33}
\]
These multiplicities follow in a basis triangularizing \(A\) and retain
coincident eigenvalues.  Since \(P_-\) commutes with \(G_M\)
and \(A^{(k)}\), image and kernel are orthogonal in that Gram
and invariant under the generator.  The full Hermitian pencil is
their orthogonal sum; its departure is \(\eta_M^-\) plus the
nonnegative complementary departure.  For \(k\ge2\), symmetric
and alternating projections are orthogonal distinct character
projections.  At \(k=1\) they coincide, so their two complements
must not be subtracted simultaneously from the identity.

\subsection{The original limit, its retained graph fibre, and determinant loss}

JD.8--9 proves \(\bigcup_{M\ge k(d-1)}\mathcal B_M\) dense in \(\mathcal H_k\)
and \(G_M\downarrow0\), by polynomial density for the original
exponentially decaying product measure.  Applying the bounded
permutation projector and (EA.13) yields
\[
 G_M^-\downarrow0,\quad \det G_M^-\to0,\quad
 R_MI_-\to0,\quad L_-J^{(k)}R_MI_-=1.
 \tag{EA.34}
\]
The latter jets remain exactly the identity at every finite cutoff.
Let \(\mathcal H_k^-\) be the alternating Hilbert subspace.  The
actual alternating polynomial source and its jet are
\[
 \mathcal D_{h,k}^-=\mathcal T_h^{(k)}
                (P_-\C[s_1,\ldots,s_k])\subset\mathcal H_k^-,
 \qquad J_-:\mathcal D_{h,k}^-\to\C^{N_-},\quad
 F\mapsto L_-J^{(k)}F.
\]
The graph closure in \(\mathcal H_k^-\times\C^{N_-}\) is precisely
\[
 \overline{\operatorname{Graph}J_-}=\mathcal H_k^-\oplus\C^{N_-},
 \quad
 \frac{\overline{\operatorname{Graph}J_-}}
                  {\mathcal H_k^-\oplus\{0\}}\simeq\C^{N_-},
 \quad [(F,u)]\mapsto u,\quad u\mapsto[(0,u)].
 \tag{EA.35}
\]
Indeed the dense alternating boundaries put \((F,0)\) in the
closure, while \((R_MI_-u,u)\to(0,u)\) puts every vertical
vector there.  Their sums exhaust the ambient product.  The displayed
quotient maps are inverse by subtraction of pairs.  Equally the exact
sequence with kernel \(u\mapsto(0,u)\), projection
\((F,u)\mapsto F\), and section \(F\mapsto(F,0)\) is split.
The limit map here is explicitly
\(\lim_M R_MI_-:\C^{N_-}\to\mathcal H_k^-\), \(u\mapsto0\),
not a continuous extension of the jet map.  Under its one-support lift it takes supported vectors
to the supported Hilbert zero, and external absence to external
absence.  Thus (EA.35) gives the exact surviving finite jet fibre.

Set \(\pi_M^- =\det G_{M+1}^-/\det G_M^-\).
The nonzero eigenvalues \(t_j\) of
\((G_M^-)^{-1/2}\Delta_M^-(G_M^-)^{-1/2}\) lie in \((0,1)\),
number \(r_M^-\), and satisfy
\(\tau_M^-=\sum_jt_j\), \(\pi_M^-=\prod_j(1-t_j)\).
Product expansion by induction and the arithmetic--geometric mean
inequality give
\[
 1-\pi_M^-\le\tau_M^-\le
 r_M^-\bigl(1-(\pi_M^-)^{1/r_M^-}\bigr),\qquad
 (\mathcal D_k^-)^2\le\tau_M^-\chi_M^-
 \quad(r_M^->0).
 \tag{EA.36}
\]
For \(r_M^-=0\), instead \(\Delta_M^-=W_M^-=0\) and
\(\pi_M^-=1\); there is no reciprocal exponent.  For integers
\(k(d-1)\le a<b\) and \(\mathcal D_k^->0\),
\[
 (\mathcal D_k^-)^2\sum_{M=a}^{b-1}(\chi_M^-)^{-1}
 \le\sum_{M=a}^{b-1}\tau_M^-
 \le\log\frac{\det G_a^-}{\det G_b^-}.
 \tag{EA.37}
\]
Every denominator is positive by (EA.30).  The second inequality
sums the exact identity
\(-\log(1-t)-t=\int_0^t u/(1-u)\,du\ge0\) and telescopes.
It also holds when \(\mathcal D_k^-=0\), without forming any
quotient.  Finally
\[
 G_a^-=\sum_{M=a}^{\infty}(Y_M^-)^*Y_M^-,\qquad
 \sum_{M=a}^{\infty}\|Y_M^-\|_{\rm HS}^2=\operatorname{Tr}G_a^-,
 \qquad \sum_{M=a}^{\infty}\tau_M^-=\infty.
 \tag{EA.38}
\]
The first two assertions follow from finite telescoping and (EA.34).
If the last series were finite, eventually all \(t_j\le\tau_M^-\le1/2\)
and \(-\log(1-t_j)\le2t_j\).  The determinant logarithm
would then remain bounded, contradicting \(\det G_M^-\to0\).
Thus determinant-line cancellation of the weight form does not cancel
the accumulated relative Gram loss.  Both facts concern the same
source sequence, through the explicit maps above.

\paragraph{Verification and scope.}
The proofs in this section are written finite algebra and source-map
arguments, not a new Lean formalization.  An independent exact checker
tests subset moments, factorial minors, nilpotent transitions, orbit
intertwining, commuting-Gram inverses, degree layers, and all small
proper-degree capacity cases.  Those checks supplement the displayed
proofs; they do not certify an arithmetic upper estimate or an RH
endpoint.  Degree \(k>d\) has zero finite exterior packet, while
degree zero in (EA.20) is only the scalar exterior-algebra convention.
